% ====================================================================== 
% CHAPTER: SYSTEM DESIGN (extracted from methodology Ch3 sec 3.8)
% ======================================================================
\chapter{System Design}
\label{ch:architecture}

\section{Overview}

The pipeline processes paired RGB and thermal-IR streams through a four-stage cascade (Figure~\ref{fig:pipeline}). The production choice for each component is justified by the ablations of Chapter~\ref{ch:hitl}.

\begin{enumerate}
    \item \textbf{RGB detector}: a YOLO26-class single-stage detector trained on the $172{,}022$-frame composite RGB dataset of Table~\ref{tab:ds_rgb_components} (drone positives, plus bird confuser negatives from the AirBird, FBD-SV, and WosDetC sources, and surveillance, traffic, and driving frames as background negatives); the production detector (\texttt{ft4}) adds the SelCom and $300$-hard-negative confuser fine-tune of Section~\ref{sec:selcom_finetune}. It is strong on drones, and residually hallucinates on small bird-on-sky frames.
    \item \textbf{IR detector}: a YOLO26-class detector trained on iteratively curated thermal imagery (Section~\ref{sec:hitl_method}), with bird, airplane, and helicopter frames as hard negatives. Strong and selective on in-distribution thermal; on the OOD thermal-confuser dataset its residual fire is airplane-dominated, and on the paired benchmarks it is the stronger single-modality detector on Svanstr\"om and the weaker one on Anti-UAV, the asymmetry the trust classifier exploits. Also operates as a \emph{grayscale-RGB fallback} when no thermal camera is available (Section~\ref{sec:ir_grayscale_mode}).
    \item \textbf{Trust classifier}: an XGBoost model that routes each frame to the trustworthy modality (\texttt{trust\_rgb}, \texttt{trust\_ir}, or \texttt{trust\_both}) and always routes, never abstaining; a fourth class, \texttt{reject\_both}, is dropped from the shipped router and kept only as an ablation (Section~\ref{sec:classifier_results}).
    \item \textbf{Confuser filter}: in production, an RGB filter that reuses the detector's own internal features, with a thermal-native IR counterpart (Section~\ref{sec:distill_verifier}). Because it reuses features already computed for each detection rather than re-processing crops, it is cheap enough to run \emph{per frame}. It supersedes an earlier patch filter, a separate CNN over crop pixels that was expensive enough to require alert-gating (Section~\ref{sec:patch_verifier_arch}).
\end{enumerate}

A temporal smoother sits on top of the cascade: it raises an alert only when the per-frame decision has persisted across a configurable sliding window. In the predecessor design the smoother also gated the (expensive) patch filter. The per-frame \texttt{mlp\_v5\_v4} filter removes that coupling, so the alert-gate analysis of Section~\ref{sec:alert_gate} now describes the superseded patch path.

\begin{figure}[htbp]
    \centering
    \includegraphics[width=\textwidth]{figures/fig_pipeline.pdf}
    \caption{Multi-stage drone detection pipeline in the production composition: detectors, then the confuser filter, then the trust classifier. The order itself was ablated (Section~\ref{sec:pipeline_ablation}): filter-first ships because it edges classifier-first on Svanstr\"om ($F1=0.946$ vs $0.931$), and the two are within about a point on the paired benchmarks. The production confuser filter is the distilled \texttt{mlp\_v5\_v4} (RGB) with its paired thermal-native IR head, which runs \textbf{per frame} because it reuses the detectors' already-computed internal features rather than re-processing crops (Section~\ref{sec:distill_verifier}). It supersedes the MobileNetV3 patch filter (grayed), which was expensive enough that it had to be \emph{alert-gated}: consulted only when the temporal smoother was about to fire (Section~\ref{sec:alert_gate}).}
    \label{fig:pipeline}
\end{figure}

Table~\ref{tab:model_overview} summarizes each production model, what it was trained on, and its held-out test result; the trust router and the two confuser filters each classify well in isolation (Figures~\ref{fig:cm_router} and~\ref{fig:cm_filters}), before any pipeline composition.

\begin{table}[htbp]
    \centering
    \caption{Production models at a glance: training data and held-out test result. Detectors report detection $F1$; the trust router and confuser filters report held-out classification (Figures~\ref{fig:cm_router}, \ref{fig:cm_filters}). The confuser filters' drone recall is reported \emph{end-to-end} (ground-truth drones kept through the detector \emph{and} the filter), so it reads below the per-detection retention in Figure~\ref{fig:cm_filters}, which conditions on detections the detector already made.}
    \label{tab:model_overview}
    \footnotesize
    \begin{tabular}{@{}p{2.9cm} l p{3.9cm} p{4.4cm}@{}}
        \toprule
        Model & Type & Trained on & Held-out test result \\
        \midrule
        RGB detector (\texttt{ft4}) & YOLO26n & $172{,}022$-frame RGB dataset $+$ SelCom/confuser fine-tune & \texttt{rgb\_dataset} test $F1$ $0.926$ \\
        \addlinespace
        IR detector (\texttt{v3b}) & YOLO26n & $129{,}130$-frame thermal dataset (6-revision HITL) & \texttt{ir\_dset} test $F1$ $0.967$ \\
        \addlinespace
        Trust router (\texttt{robust8-nr}) & XGBoost & $37{,}488$ routing-labeled detections & held-out fold ($11{,}318$ dets): macro-$F1$ $0.941$, accuracy $0.981$ \\
        \addlinespace
        RGB filter (\texttt{mlp\_v5\_v4}) & MLP & $33{,}272$ \texttt{ft4} features ($19{,}334$ drone / $13{,}938$ confuser) & \texttt{rgb\_dataset} test drone recall $0.887$; held-out bird $30/230$ kept \\
        \addlinespace
        IR filter (\texttt{mlp\_aligned\_\allowbreak thermalonly}) & MLP & $10{,}157$ \texttt{v3b} features ($8{,}112$ drone / $2{,}045$ confuser) & \texttt{ir\_dset} test drone recall $0.928$; CBAM FP $48\to6$ \\
        \bottomrule
    \end{tabular}
\end{table}
% [source: RGB det F1 0.926 = conf_sweep_results.json rgb_dataset_test peak (sec:design_rationale); IR det F1 0.967 = tab:ir_evolution; router macro-F1 0.941/acc 0.981 = thesis_eval/eval_router_heldout.py -> results/per_model_heldout/router_heldout.json; RGB filter recall 0.887 = sec:distill_verifier, bird 30/230 = eval/eval_birdtest_heldout.py; IR filter recall 0.928 + CBAM 48->6 = tab:ir_aligned; trained-on counts = model cards + _v5_balanced_v4/training_data.npz + fusion_dataset_full56]

The downstream stages contribute the majority of confuser suppression without retraining the detectors; Figure~\ref{fig:confuser_problem} shows the failure mode this addresses and the per-detection separation that fixes it. Three architectural choices make this possible: \emph{deferred suppression} at the detectors (the detectors hold standard operating points and are never made stricter to fight confusers; precision is owned downstream, Section~\ref{sec:design_rationale}), trust-aware fusion (Section~\ref{sec:trust_fusion}), and a downstream per-frame confuser filter. The production filter is a confidence-gated veto: a detection survives only if the MLP assigns it $P(\text{drone})$ at or above the production threshold (RGB $0.25$, IR-thermal $0.05$). A fail-open filter variant that releases uncertain detections was evaluated and rejected, because the detections it releases on cluttered surfaces are exactly the false positives the filter exists to remove.
% [source: docs/analysis/2026-06-01_mlp_v5_recall_drop_mri.md (fail-open conditional: recovers rgb_dataset_test recall but craters svanstrom precision 0.887->0.631)]

\begin{figure}[htbp]
    \centering
    \includegraphics[width=0.85\textwidth]{fig_confuser_panel}
    \caption{The confuser problem, and why a downstream filter solves it. \emph{Left}: the production RGB detector (\texttt{ft4}) fires on a bird amid its flock (red box, detector confidence $0.46$). \emph{Right}: it fires on a drone (confidence $0.85$). The detector cannot tell them apart, but the feature-reuse filter reads the detector's own internal features and assigns $P(\text{drone})=0.00$ to the bird (vetoed) and $0.96$ to the drone (kept). This separation, not a better detector, is what the cascade provides.}
    \label{fig:confuser_problem}
\end{figure}

\paragraph{Two-filter fusion (filter-first).} With both the RGB (\texttt{mlp\_v5\_v4}) and IR filters running per-frame, the production composition is \emph{filter-first}: each modality's confuser filter screens its detections, then the trust classifier (Section~\ref{sec:trust_fusion}) routes among the survivors. \texttt{trust\_rgb} keeps the detection only if the RGB filter passed; \texttt{trust\_ir} keeps it only if the IR filter passed; \texttt{trust\_both} is resolved per-modality (recall-first: the detection survives if either trusted filter passed). The shipped router never abstains; the reject-capable ablation adds a \texttt{reject\_both} outcome that drops the frame outright. The IR filter is a single thermal-native CBAM net (\texttt{mlp\_aligned\_thermalonly}, $P(\text{drone})\geq 0.05$; Section~\ref{sec:ir_xmodal_verifier}). The composition order (filter$\to$classifier vs classifier$\to$filter) is itself ablated in Section~\ref{sec:pipeline_ablation}.
% [source: two-filter fusion rule = gui/pyside_engine.py (trust\_both resolved recall-first: detection survives if either trusted filter passes)]

\section{Design Rationale: Deferred Suppression}
\label{sec:design_rationale}

The cascade's central design principle is asymmetric recoverability: a confuser false positive raised by the detector can be vetoed downstream, but a drone the detector missed cannot be recovered by any later stage. Suppression is therefore never the detector's job. This does not mean the detectors run loose. Their floors are standard operating points (RGB $0.25$, IR $0.40$ in deployment; the canonical evaluation caches both at $0.25$), and the confidence sweep shows those are normal optima rather than opened gates: Anti-UAV bare $F1$ is flat ($0.959$--$0.963$) across floors $0.05$--$0.5$, and the in-distribution RGB test split peaks at exactly $0.25$ ($F1=0.926$). What the principle forbids is moving the floor \emph{up}, or retraining the detector into suppression, when confusers appear. Both detectors are trained with confuser negatives (Section~\ref{sec:rgb_three_variants}, Section~\ref{sec:ir_detector}), but training-time mining does not eliminate all categories. Bird hallucination at the RGB detector is irreducible at small drone sizes (Figure~\ref{fig:dataset_montage}): the three-stance comparison shows that the stance that suppresses birds (the aggressively-retrained detector, bird fire $94.4\% \to 3.4\%$) collapses Svanstr\"om drone recall $0.961 \to 0.306$ (Table~\ref{tab:rgb_comparison}, Section~\ref{sec:rgb_results}). The cascade is built around that residual: the downstream stages catch what training cannot, at a cost and benefit quantified stage-by-stage in Section~\ref{sec:pipeline_ablation}. The same deferral is what licenses the low-confidence operating mode (Section~\ref{sec:lowconf_mode}): on a recall-starved OOD surface the floor can drop to $0.05$ \emph{because} precision ownership has already moved downstream (SelCom $0.591$ bare at the default floor $\to$ $0.692$ with the filter at floor $0.05$).
% [source: thesis_eval/results/conf_sweep/conf_sweep_results.json (antiuav rgb-slice bare f1 0.9592@0.05..0.9631@0.25..0.9603@0.5; rgb_dataset_test bare 0.9259@0.25 peak; selcom_val bare 0.5911@0.25, +filt 0.692@0.05)]
% [source: ledger=retrainedv2-recall-collapse; eval=rgb_svan_retrainedv2; cache=eval/results/_failure_diagnosis/svanstrom_1280_by_category.csv; config=svan_iop_1280]

\section{Trust-Aware Fusion}
\label{sec:trust_fusion}

The fusion stage is a learned per-frame decision rather than a fixed vote. The production router outputs one of three labels and always routes:

\begin{itemize}
    \item \textbf{\texttt{trust\_rgb}}, only RGB is reliable for this frame.
    \item \textbf{\texttt{trust\_ir}}, only IR is reliable for this frame.
    \item \textbf{\texttt{trust\_both}}, both modalities agree and are trustworthy.
\end{itemize}

The classifier is trained with a fourth label, \texttt{reject\_both} (neither modality's detection is trustworthy, likely a confuser), but the shipped router drops it (leaving false-positive rejection to the per-frame filter) and keeps the reject-capable variant only as an ablation (Section~\ref{sec:classifier_results}).

The routing formulation captures \emph{when} each modality is reliable. The asymmetry the classifier exploits is that \emph{the two detectors hallucinate on different scenes}: visual confusers fool the RGB detector more reliably; thermally degraded conditions fool the IR detector more reliably. The router does not read the scene to detect this. It sees \emph{detection evidence only}: its eight features are the per-modality maximum detection confidence, the mean RGB confidence, and the best detection's log box area and aspect ratio for each modality, plus the binary grayscale-regime flag. A frame on which one detector emits low-confidence, oddly-shaped boxes while the other emits a single confident, well-proportioned one carries the failure-profile signal in the detections themselves; scene statistics appear only in the hand-engineered \texttt{sa32} comparison classifier. A per-frame classifier can pick the modality whose failure profile is least active; a fixed-weight rule cannot. The per-surface evidence, including the modality reversal between Svanstr\"om (IR stronger) and Anti-UAV (RGB stronger) that makes routing pay, is in Section~\ref{sec:pipeline_ablation}.

The router's eight features were screened by the leakage statistic of Section~\ref{sec:model_mri}, which drops any feature that discriminates \emph{which scene} a frame came from rather than \emph{whether the detections are trustworthy}, so the router cannot pass its evaluation by memorizing surface identity. The strongest hand-built comparison, a 32-feature classifier, leaks exactly those scene statistics, so its in-domain lead is read as an upper bound rather than a deployable target.
% [source: classifier/train_routing_robust.py ROBUST6 + rgb_mean_conf + is_grayscale; knowledge/models.csv robust8 ("8 free conf+geometry feats")]

\begin{figure}[htbp]
    \centering
    \includegraphics[width=0.52\textwidth]{cm_router}
    \caption{Trust router (\texttt{robust8-nr}) held-out classification, on a grouped split (by sequence) of its dataset, shipped model not retrained: macro-$F1$ $0.941$, accuracy $0.981$. The dominant off-diagonal (\texttt{rgb}$\to$\texttt{both}) is recall-safe.}
    % [source: thesis_eval/eval_router_heldout.py -> thesis_eval/results/per_model_heldout/router_heldout.json]
    \label{fig:cm_router}
\end{figure}

\section{Temporal Smoother}
\label{sec:temporal}

The temporal smoother is an $N$-of-$M$ sliding-window aggregator with configurable cooldown. An alert fires only when the trust classifier's per-frame decision has been positive on at least $N$ of the last $M$ frames, absorbing isolated firings characteristic of birds and aircraft transiting the field of view. Its measured contribution, at the segment grain on the real-video clips, is in Section~\ref{sec:temporal_results}.

The same temporal state drives a \emph{ROI fallback} in the PySide engine (\texttt{\_run\_with\_roi}). When the full-frame YOLO returns nothing and a recent track is still alive, YOLO is re-run on an expanded crop around the previous box at reduced confidence. This is a GUI-only feature, and the frame-mode evaluation harness does not use it; it accounts for part of the GUI-versus-eval performance gap on small distant drones. The GUI maintains tracks and renders a path image, but does not yet classify on track-level features (flight dynamics); track-level classification is a future-work direction.

\begin{figure}[htbp]
    \centering
    \begin{subfigure}{0.49\textwidth}\includegraphics[width=\textwidth]{gui_grayscale}\caption{Grayscale fusion}\end{subfigure}
    \hfill
    \begin{subfigure}{0.49\textwidth}\includegraphics[width=\textwidth]{gui_paired_thermal}\caption{Paired RGB$+$thermal fusion}\end{subfigure}
    \caption{PySide6 operator GUI (\texttt{TALOS}) on real-video footage. The per-frame trust decision (\textsc{trust both}/RGB/IR), the confuser-filter verdict, the alert counters, and the temporal-window progress are surfaced for the operator: (a) grayscale fusion and (b) paired RGB$+$thermal fusion.}
    \label{fig:pyside_gui}
\end{figure}

\section{Resolution Dependency}
\label{sec:resolution_arch}

Resolution is the single largest hyperparameter affecting small-drone recall. Svanstr\"om's native resolution is $640\times 512$; a dedicated sweep evaluating both RGB variants at both sizes under one harness (stride-7 sample, $n=4{,}102$, IoP@0.5) isolates the resolution effect within each model (Figure~\ref{fig:resolution}): doubling the inference size buys both models roughly the same recall increment ($-28$~pp baseline, $-25$~pp \texttt{retrained\_v2} at 640), while the absolute level is set by the training stance, with the confuser-suppressing \texttt{retrained\_v2} unusable on Svanstr\"om at either size. On the full dataset the baseline's canonical 1280 figure is $R=0.961$. Anti-UAV imagery ($1920\times 1080$) saturates at \texttt{imgsz=640} because drones already span enough pixels. All Svanstr\"om RGB evaluation uses \texttt{imgsz=1280} (Section~\ref{sec:canonical_configs}). SelCom CCTV is a third regime: clutter is the binding constraint, and the imgsz sweep identifies $\mathtt{imgsz}=960$ as the F1-optimal setting for that camera's deployment operating point.
% [source: eval/results/svan_resolution_sweep.json (baseline@640 R 0.6838 / @1280 0.9641; retrained_v2@640 0.0699 / @1280 0.3234; stride 7, n=4102, conf 0.25, IoP@0.5); run=eval/svan_resolution_sweep.py; full-dataset baseline 0.961 = evals.csv rgb_svan_baseline]
% [source: ledger=selcom-imgsz-win; eval=rgb_svan_baseline; run=eval/diagnose_failures.py; cache=eval/results/_failure_diagnosis/svanstrom_1280_by_category.csv; config=svan_iop_1280]

\begin{figure}[htbp]
    \centering
    \includegraphics[width=0.85\textwidth]{fig6_6_resolution}
    \caption{Inference-resolution dependency of Svanstr\"om drone recall, isolated \emph{within} each model (one harness: stride-7 sample, $n=4{,}102$, IoP@0.5, conf $0.25$). Doubling \texttt{imgsz} buys both variants a similar recall increment ($+28$~pp baseline, $+25$~pp \texttt{retrained\_v2}); the absolute level is set by the training stance. The confuser-suppressing \texttt{retrained\_v2} is unusable on Svanstr\"om at either size.}
    % [source: ledger=retrainedv2-recall-collapse; eval=rgb_svan_baseline; run=eval/diagnose_failures.py; cache=eval/results/_failure_diagnosis/svanstrom_1280_by_category.csv; config=svan_iop_1280]
    \label{fig:resolution}
\end{figure}

% tab:resolution removed (duplicate of fig:resolution; exact values retained in the figure source comment + svan_resolution_sweep.json).

\section{RGB Drone Detector}
\label{sec:rgb_detector}

The RGB detector is a YOLO26n (nano) single-class detector trained on the composite dataset of Section~\ref{sec:ds_rgb_corpus}, which includes a drone-vs-bird subset as deliberate confuser negatives. Svanstr\"om imagery is held out from training.

\subsection{Three Training Stances}
\label{sec:rgb_three_variants}

Three RGB-detector variants explore how much additional confuser-specific data to add on top of the baseline dataset; they occupy three operating points on the recall--hallucination axis, and their full comparison is Table~\ref{tab:rgb_comparison} (Section~\ref{sec:rgb_results}).

\begin{enumerate}
    \item \textbf{Baseline.} Trained on the composite drone dataset. On Svanstr\"om@1280 the bare detector, no downstream stage applied, reaches $P=0.940, R=0.961$ on drones, with a 94.4\% bird-frame fire rate.
    % [source: knowledge/evals.csv rgb_svan_baseline (S1 RGB-alone, full 28,710); eval/results/_failure_diagnosis/svanstrom_1280_by_category.csv (P 1248/1327, R 1248/1299, bird 807/856 frames)]
    \item \textbf{Hard-negative training.} Same backbone, retrained with the Svanstr\"om bird/\allowbreak airplane/\allowbreak helicopter confuser splits supplied as additional negatives. Helicopter fire drops $66.2\% \to 41.9\%$ and airplane $74.6\% \to 64.7\%$, while bird fire is essentially unchanged ($94.2\%$): hard-negative mining is class-specific.
    \item \textbf{Aggressive retraining.} With heavy confuser augmentation, bird fire collapses to 3.4\%, and Svanstr\"om drone recall collapses with it, to $R=0.306$. The lesson that motivated the per-detection confuser filter: at small drone scales, bird suppression and drone recall cannot both be had at the detector. (The trust router is needed regardless; its job is modality choice, not confuser rejection.)
\end{enumerate}
% [source: ledger=retrainedv2-recall-collapse; eval=rgb_svan_baseline,rgb_svan_hardneg,rgb_svan_retrainedv2; cache=eval/results/_failure_diagnosis/svanstrom_1280_by_category.csv; config=svan_iop_1280]

The production RGB detector is a fourth stance built on the lessons of the first three: the SelCom fine-tune line (Section~\ref{sec:selcom_finetune}) plus a \emph{regression-gated} confuser injection (300 hard negatives at \texttt{freeze}$=$15; Section~\ref{sec:training_recipes}) that cut confuser hallucination 16~pp without recall loss.

\subsection{SelCom CCTV Fine-Tune}
\label{sec:selcom_finetune}

A separate fine-tune line targets the deployment partner's CCTV stream (Section~\ref{sec:ds_selcom}). The fine-tune is on top of the \emph{baseline} weights (not \texttt{retrained\_v2}). Two mixed-training fine-tunes (80\% general RGB + 20\% selcom; pure-selcom val split, seed 0) were trained at 640 and 1280; the 1280 variant doubles SelCom recall \emph{and} raises precision over its 640 sibling (Table~\ref{tab:selcom}, Section~\ref{sec:rgb_results}), the clearest single demonstration that resolution, not architecture, is the binding constraint on small-drone CCTV. A preprocessing sweep established that no OpenCV variant (denoise, CLAHE, sharpen, contrast) improves detection independently of fine-tuning and \texttt{imgsz}; the imgsz sweep identifies 960 as the F1-optimal deployment point for this camera. The regression on the standard RGB dataset is within $-0.008$ F1 of baseline.
% [source: ledger=selcom-imgsz-win; eval=selcom_ft2_1280; run=RGB model/finetune_selcom.py; cache=runs/rgb_finetune_eval/Yolo26n_selcom_mixed_ft2_1280/comparison.json; config=selcom_iop_1280]

\section{IR Drone Detector}
\label{sec:ir_detector}

The IR detector uses the same YOLO26n architecture finetuned on iteratively curated thermal data via the HITL engine (Section~\ref{sec:hitl_method}). Bird, airplane, and helicopter thermal frames are added as hard negatives across revisions. The production weights (\texttt{v3b}) are a 2-epoch corrective fine-tune on the \texttt{Final} checkpoint addressing specific deployment-time failure modes. The full evolution across six numbered revisions (including the V5 regression and its recovery) is the case study of Section~\ref{sec:ir_evolution}; on the paired benchmarks the IR detector is the stronger single-modality detector on Svanstr\"om and the weaker one on Anti-UAV (Section~\ref{sec:pipeline_ablation}), which is precisely the asymmetry the trust classifier exploits.

\subsection{Grayscale-RGB Operating Mode}
\label{sec:ir_grayscale_mode}

The deployed PySide GUI exposes a \emph{grayscale-RGB} operating mode (\texttt{\_infer\_\allowbreak grayscale}). When no thermal camera is available, the GUI converts the RGB stream to grayscale via \texttt{cvtColor(\allowbreak BGR2GRAY)}, replicates to three channels, and feeds the result to the IR detector at \texttt{imgsz=640}. This was built as a conservative hardware fallback; its measured behavior (competitive with the dedicated RGB detector on bird-cluttered scenes, a few points behind it overall) turned out to be one of the thesis's findings (Section~\ref{sec:grayscale}).

\subsection{Cross-Modal Feature Alignment: a Thermal-Native Confuser Filter}
\label{sec:ir_xmodal_verifier}

The IR branch carries a per-detection confuser filter of the same kind as the RGB branch: a thermal-native net trained directly on thermal features. A Model MRI measurement (evidence figures in Section~\ref{sec:mri_findings}) establishes that the signal it needs is already present in the detector:

\emph{The confuser-rejection signal already lives inside the IR detector.} A Model MRI scan (the feature-space instrument of Section~\ref{sec:model_mri}) of the \texttt{v3b} fused \texttt{p3}+\texttt{p5} features, over $14{,}697$ drone and $1{,}386$ confuser detections, finds the two classes almost linearly separable (Table~\ref{tab:ir_mri_sep}): a single linear discriminant splits them at $0.981$ train accuracy. The IR feature space is the thermal twin of the RGB filter's, readable by the same kind of lightweight MLP with no detector retraining. The IR detector is also, unlike the RGB detector, already clean on near-domain confusers: on its Model-MRI thermal-confuser dataset (Section~\ref{sec:model_mri}, the IR mining set) it fires on only $1.8\%$ of images, against the RGB detector's $30.4\%$ frame-fire on the \texttt{rgb\_confusers\_merged} out-of-distribution dataset (Section~\ref{sec:pipeline_confusers}), so a thermal confuser filter is not strictly required. It earns its place on the hardest out-of-distribution thermal-aerial confusers: on the held-out CBAM \texttt{valid} split (Section~\ref{sec:ds_ir_confusers}; the TEST split disjoint from the CBAM \texttt{train} drones the head is trained on), the thermal-native head recovers drone recall to $R=0.967$ while cutting false positives to $6$ (bare $48$), at no recall cost (Section~\ref{sec:verifier_results}).
% [source: ledger=ir-features-separable; run=mri/cli.py; cache=mri/results/ir_v3b_report (stats.json); rgb 0.304 frame-fire = tier1 rgb_confuser bare; CBAM R 0.967/FP 6 (bare 48) = mlp_aligned_thermalonly held-out CBAM valid, docs/analysis/2026-06-18_filter_provenance_train_heldout.md §2.4; config=ir_v3b]

\begin{table}[htbp]
    \centering
    \caption{Model MRI of the \texttt{v3b} IR detector's fused \texttt{p3}+\texttt{p5} features ($14{,}697$ drone / $1{,}386$ confuser detections): the confuser-rejection signal is already linearly present in the detector's own features. The projected columns are in-pool $5$-fold estimates of what a reader of those features could recover.}
    \label{tab:ir_mri_sep}
    \begin{tabular}{lc}
        \toprule
        MRI signal & value \\
        \midrule
        LDA train separability (drone vs confuser) & $0.981$ \\
        Max ANOVA $F$ (strongest single neuron) & $5{,}370$ (\texttt{p5}) \\
        Median ANOVA $F$ & $256$ \\
        Raw detector hallucination rate & $1.8\%$/img \\
        Projected FP cut ($\leq$10\% recall cost) & $89\%$ \\
        Projected recall retention & $99.7\%$ \\
        \bottomrule
    \end{tabular}
\end{table}
% [source: mri/results/ir_v3b_report/stats.json (lda 0.9813, max_anova_F 5370.1, median 255.7, raw_halluc 0.01765, fp_reduction 0.8896, recall_retention 0.99687); n=14697/1386]

Production therefore ships a single $517$-D thermal-native head (\texttt{mlp\_aligned\_\allowbreak thermalonly}, a CBAM-trained net), trained on $8{,}112$ thermal drones and $2{,}045$ \texttt{IR\_confusers/train} confusers (Svanstr\"om-IR, Anti-UAV-val-IR, \texttt{ir\_dset\_final}-train thermal drones plus the CBAM \texttt{train} drones, GT-aware). The thermal-native head holds drone recall on held-out thermal surfaces (antiuav\_ir $0.937$, ir\_video $0.971$, svanstrom\_ir $0.966$); the only cost is \texttt{ir\_dset\_final} $0.965 \to 0.928$ ($-3.7$~pp), on genuinely airplane-like drones that no threshold separates. It cuts CBAM confuser false positives from $48$ to $6$ (held-out CBAM \texttt{valid} split). The thermal head is trained on \texttt{IR\_confusers/train} confusers plus Svanstr\"om-IR / Anti-UAV-val / \texttt{ir\_dset\_final}-train thermal drones, and is evaluated on the held-out CBAM \texttt{valid} split and the \texttt{IR\_confusers} val/test re-mine (TEST split, disjoint from training); the held-out results are in Section~\ref{sec:verifier_results}.
% [source: docs/analysis/2026-06-18_filter_provenance_train_heldout.md §2.4 (thermal-only mlp_aligned_thermalonly @0.05: held-out CBAM valid R 0.967 / FP 6 vs bare 48; antiuav_ir 0.937, ir_video 0.971, svanstrom_ir 0.966 held; ir_dset_final 0.965->0.928 -3.7pp; IR_confusers val/test held-out 90->22); knowledge/models.csv mlp_aligned_thermalonly]
% [source: docs/analysis/2026-06-18_thesis_integration_checklist.md §2 (CBAM tab:ir_aligned RESOLVED R 0.967/FP 6); ledger=filter-stack-v4-thermalonly-promoted]

\section{Trust Classifier}
\label{sec:trust_classifier}

\subsection{Feature Set}

The shipped trust classifier consumes per-frame features computed from the two detectors' raw outputs alone (detection flags, confidences, box geometry, cross-modal agreement); it does not read the source frames. The shipped router uses eight features and always routes (the reject-capable variant and the ship rationale are in Section~\ref{sec:classifier_results}). The hand-built comparison set adds a 32-feature classifier and two 40-feature variants that also read source-frame scene statistics, which is the family the leakage analysis of Section~\ref{sec:feature_selection} flags and drops. The feature families:

\begin{itemize}
    \item \textbf{Detector signals}: per-modality binary detection flags, top-$k$ confidence scores, detection counts.
    \item \textbf{Scene features}: image intensity statistics, edge density, frequency-domain summaries; intended to detect ``confuser-rich'' scenes vs ``clean sky'' scenes; as Section~\ref{sec:feature_selection} shows, also the family that leaks scene identity.
    \item \textbf{Target features}: bounding-box geometry (aspect ratio, normalized area, position), confidence distribution per modality.
    \item \textbf{Agreement features}: cross-modal IoU, center distance, confidence ratio.
\end{itemize}

The label space is the 4-way trust decision of Section~\ref{sec:trust_fusion}. The model is an XGBoost~\cite{chen2016xgboost} gradient-boosted classifier. Training uses sequence-stratified group shuffling so that frames from a given sequence are not split across train and test. The Svanstr\"om overlap this leaves at the cascade level is flagged in Section~\ref{sec:svanstrom_audit}.

\subsection{Hand-Engineered Variants (the comparison set)}
\label{sec:classifier_variants}

Three hand-engineered variants preceded the statistical selection: \texttt{sa32} (\texttt{scene\_aware\_\allowbreak v3more\_\allowbreak 32feat}; the strongest, and the working production pick through the real-video comparison, though it leaks scene statistics so its lead is not trusted to generalize), \texttt{fusion\_no\_fn\_v1.1} (the most conservative on OOD confusers; rejects most drone TPs on operational video; retained as an open-world fallback where missed drones are operationally tolerable), and \texttt{control\_\allowbreak v3more\_\allowbreak 40feat} (deprecated; no surface favors it). Their quantitative comparison, and the surface reversal between Svanstr\"om and real video that it revealed, is in Section~\ref{sec:classifier_results}.

The early \texttt{retrained\_v2\_32feat} classifier is superseded because the production RGB detector is no longer \texttt{retrained\_v2}; a training-time/production confidence-distribution mismatch accounts for most of its gap. A trust classifier is not portable across RGB detector swaps without re-checking calibration.

\paragraph{Scene-fingerprint over-fitting.}
Under sequence-stratified splitting, per-clip image statistics act as \emph{scene fingerprints}: the classifier learns which clip a frame came from rather than whether a drone is present, inflating in-distribution scores while degrading transfer. Dropping these features (the \texttt{lean13}$\to$\texttt{lean10} pruning) recovers $18$--$26$~pp drone $F1$ on held-out clips. Section~\ref{sec:feature_selection} turns this hindsight observation into a statistic computed before training.
% [source: ledger=scene-fingerprint-overfit; eval=own_lean10_yt,own_lean13_yt; config=clf_own_holdout]

\subsection{Statistical Feature Selection for the Router}
\label{sec:feature_selection}

The variants above were assembled by hand; this subsection replaces that procedure with a statistical one. Two practical pressures sharpen the motivating question: \texttt{sa32} was trained on detections from a superseded detector (its inputs no longer match the production \texttt{ft4}/\texttt{v3b}), and several of its features require an expensive per-frame OpenCV pass.

\paragraph{Method.}
From a cached $56$-feature matrix of $65{,}192$ frames we ran four analyses. Linear Discriminant Analysis (LDA), which projects features onto the single axis that best separates the two classes, asks whether the signal is present: a single Fisher axis separates trust-positive from reject at $98.2\%$ training accuracy (Figure~\ref{fig:fusion_lda}), so the problem is \emph{selecting} features, not a signal shortage. Principal Component Analysis (PCA) asks whether the structure is trivial: classes overlap in the top components, meaning scene-to-scene variance dominates over the drone signal. A per-feature ANOVA $F$-test and single-feature AUROC (area under the ROC curve) then rank each feature individually (Figure~\ref{fig:fusion_auroc}).
% [source: ledger=fusion-feature-leakage; run=classifier/fusion_feature_stats.py; cache=classifier/fusion_models/optimal_v1/feature_stats_ranked.csv; config=clf_own_holdout]

\paragraph{A leakage statistic.}
A scene-fingerprint feature can score as discriminative in a single pool simply because it correlates with which scene a sample came from; ANOVA and AUROC would recommend keeping it. To expose this we compute a second $F$-statistic \emph{within the drone class only} and form the ratio
\begin{equation}
\mathit{leakage\_ratio} = \frac{F_\text{domain-in-class}}{F_\text{class}},
\label{eq:leakage}
\end{equation}
which is low for a robust signal feature and high for a scene fingerprint. Equation~\ref{eq:leakage} converts the earlier hand-tuned deprecations into a number computed before any training run.
% [source: ledger=fusion-feature-leakage; run=classifier/fusion_feature_stats.py; cache=classifier/fusion_models/optimal_v1/feature_stats_ranked.csv; config=clf_own_holdout]

\begin{figure}[htbp]
    \centering
    \begin{subfigure}{0.32\textwidth}\includegraphics[width=\textwidth]{fig8_fusion_lda}\caption{}\label{fig:fusion_lda}\end{subfigure}
    \hfill
    \begin{subfigure}{0.32\textwidth}\includegraphics[width=\textwidth]{fig8_fusion_auroc}\caption{}\label{fig:fusion_auroc}\end{subfigure}
    \hfill
    \begin{subfigure}{0.32\textwidth}\includegraphics[width=\textwidth]{fig8_fusion_leakage}\caption{}\label{fig:fusion_leakage}\end{subfigure}
    \caption{Statistical anatomy of the $56$-feature fusion space. (a) Linear discriminant: the trust signal is linearly separable ($98.2\%$ train accuracy). (b) Per-feature single-feature AUROC. (c) Leakage map (Eq.~\ref{eq:leakage}): the lower-right region is robust signal (high AUROC, near-zero leakage); the upper-left is scene fingerprint (chance-level AUROC, leakage in the hundreds). Source: \texttt{classifier/fusion\_feature\_stats.py}.}
    % [source: ledger=fusion-feature-leakage; run=classifier/fusion_feature_stats.py; cache=classifier/fusion_models/optimal_v1/feature_stats_ranked.csv; config=clf_own_holdout]
    \label{fig:fusion_stats}
\end{figure}

\paragraph{What the statistic finds.}
The leakage map separates two groups cleanly (Table~\ref{tab:leakage}). Scene-statistic features score at chance on the class label yet have leakage ratios in the hundreds: they discriminate \emph{which scene}, not whether a drone is present, and are also the expensive per-frame features. The robust core (detection confidences, box geometry, and cross-modal ratios) scores AUROC above $0.9$ with leakage below $0.006$. A within-source check catches what the leakage ratio misses: \texttt{rgb\_blurriness} scores AUROC $0.872$ in the pooled set but collapses to $0.516$ within each source (a cross-clip mean-shift artefact), so the gate drops it alongside the scene fingerprints. Two caveats. Per-modality detection-presence flags top the AUROC chart but are excluded by reasoning, since the trust label is \emph{derived from} those detections. And the dataset is IR-dominant, so IR features rank above their RGB counterparts on RGB-fallback surfaces.
% [source: ledger=fusion-feature-leakage,blurriness-is-dataset-artifact; run=classifier/fusion_feature_stats.py,classifier/vet_blurriness_stats.py; cache=classifier/fusion_models/optimal_v1/feature_stats_ranked.csv; config=clf_own_holdout]

\begin{table}[htbp]
    \centering
    \caption{Leakage analysis splits the feature space. \emph{Top}: scene fingerprints, chance-level class AUROC but very high leakage; drop. \emph{Bottom}: the robust core, high AUROC, near-zero leakage; keep. Source: \texttt{feature\_stats\_ranked.csv}.}
    \label{tab:leakage}
    \begin{tabular}{llcc}
        \toprule
        feature & family & AUROC-alone & leakage ratio \\
        \midrule
        \multicolumn{4}{l}{\emph{Scene fingerprints (drop)}} \\
        \texttt{rgb\_img\_std}       & scene       & 0.502 & 349.6 \\
        \texttt{rgb\_img\_entropy}   & scene       & 0.510 & 307.4 \\
        \texttt{ir\_blurriness}      & scene       & 0.795 & 11.7 \\
        \texttt{rgb\_edge\_density}  & scene       & 0.525 & 1.75 \\
        \midrule
        \multicolumn{4}{l}{\emph{Robust core (keep)}} \\
        \texttt{conf\_sum}             & confidence  & 0.983 & 0.002 \\
        \texttt{ir\_max\_conf}         & confidence  & 0.965 & 0.002 \\
        \texttt{ir\_best\_aspect\_ratio} & geometry  & 0.952 & 0.002 \\
        \texttt{ir\_best\_log\_bbox\_area} & geometry & 0.946 & 0.005 \\
        \texttt{xmodal\_conf\_ratio}   & cross-modal & 0.905 & 0.002 \\
        \texttt{rgb\_max\_conf}        & confidence  & 0.816 & 0.004 \\
        \texttt{rgb\_best\_log\_bbox\_area} & geometry & 0.806 & 0.059 \\
        \texttt{rgb\_best\_aspect\_ratio} & geometry & 0.778 & 0.011 \\
        \bottomrule
    \end{tabular}
\end{table}

\paragraph{\texttt{robust6}.}
Applying the keep rule yields \texttt{robust6} $=$ \{\texttt{rgb\_max\_conf}, \texttt{ir\_max\_conf}, \texttt{rgb\_best\_\allowbreak log\_bbox\_area}, \texttt{ir\_best\_\allowbreak log\_bbox\_area}, \texttt{rgb\_best\_\allowbreak aspect\_ratio}, \texttt{ir\_best\_\allowbreak aspect\_ratio}\}: the six independent per-modality primitives, all free detector outputs and box geometry, all near-zero leakage. The two remaining keep-set entries in Table~\ref{tab:leakage}, \texttt{conf\_sum} and \texttt{xmodal\_conf\_ratio}, clear the same gate but are summary aggregates of those per-modality confidences, so \texttt{robust6} does not carry them separately; the table lists a representative slice of the $56$-feature pool, not the final feature set. Training data is re-mined from the current \texttt{ft4}/\texttt{v3b} detectors, fixing the drift in \texttt{sa32}'s inputs. Two ablation levels confirm the design: in-domain $F1$ is flat as the feature count grows from $6$ to $32$ (count is not the lever), but on OOD drone video the set \emph{with} fingerprint features collapses to $F1=0.262$, climbing to $0.578$ for \texttt{robust6}; dropping fingerprints roughly doubles OOD generalization. Dropping IR geometry also loses confuser rejection, so box-geometry features are load-bearing. \texttt{robust8} extends this six-feature core with two further free detector features (\texttt{rgb\_mean\_conf} and \texttt{is\_grayscale}); the shipped router uses those same eight features but drops the \texttt{reject} class (\texttt{robust8-nr}), for the reasons given in Section~\ref{sec:classifier_results}.
% [source: ledger=ft4-lean-trust-classifier; run=classifier/train_lean_ft4.py; cache=docs/analysis/2026-05-31_ft4_lean_trust_classifier.md; config=clf_own_holdout]

\paragraph{A further free lever.}
The strongest unused lever is the filter's own output: the RGB filter's drone probability separates trust-positive from reject at AUROC $0.949$ vs $0.842$ for raw confidence. It is withheld from \texttt{robust8} because coupling the filter into the classifier introduces stage dependency; it is recorded as the next free lever.
% [source: run=classifier/verifier_as_feature_stats.py; cache=docs/analysis/2026-06-02_forward_selection_study.md]

\section{The Confuser Filter}
\label{sec:distill_verifier}

The patch filter re-processes each detection crop with a network entirely separate from the detector; on unseen surfaces it abstains to \texttt{other} (on SelCom CCTV it is bit-for-bit identical to the bare detector, having never seen CCTV-scale crops). This motivates a structurally different filter.
% [source: ledger=patch-v2-neutral-selcom; eval=v5_selcom_patch; config=selcom_iop_1280]

The \emph{feature-reuse filter} replaces pixel re-processing with feature reuse: a small MLP consumes the RGB detector's own internal \texttt{p3}+\texttt{p5} activations and detection metadata, the same $517$-D per-detection embedding the Model MRI extracts (Section~\ref{sec:model_mri}), already computed for each candidate box. The production filter extends this dataset with an in-distribution bird split and a rebalanced size-by-source drone manifold, which closes a $22\%$ coverage gap where the predecessor over-vetoed real \texttt{rgb\_dataset\_test} drones. The Model MRI characterized the parent feature space on the $32{,}931$-detection dataset the shipped predecessor was trained on ($19{,}334$ drone / $13{,}597$ confuser detections mined across the RGB surfaces). A single linear discriminant separates drone from confuser at ${\approx}95\%$ train accuracy, while unsupervised PCA shows near-zero cluster separation (silhouette $0.067$): the signal is real but lives in a supervised subspace, so it needs a trained classifier rather than a distance threshold. The MRI findings, including the activation ``brain scans'' of what the filter reads, are in Section~\ref{sec:mri_findings}; the filter's measured effect on every surface is in Section~\ref{sec:verifier_results}. \texttt{mlp\_v5\_v4} is trained on \texttt{rgb\_dataset} train$+$val, \texttt{rgb\_confusers} train$+$val, pure-SelCom (the $311$ \texttt{selcom\_val} files blocklisted), and the bird.v1i $728$-image \emph{train} split, in all $33{,}272$ detections ($19{,}334$ drone / $13{,}938$ confuser, Table~\ref{tab:model_overview}); the $32{,}931$ above is that same dataset \emph{before} its $341$-confuser bird split, and is what the Model MRI characterized. It is evaluated on the held-out bird.v1i $484$-image \emph{test} split (TEST split, disjoint by name from training): it keeps only $30/230$ unseen-bird fires against the shipped filter's $91/230$, and lifts \texttt{rgb\_dataset\_test} recall $0.691 \to 0.887$. All numbers are reported on each dataset's TEST split; Svanstr\"om-RGB is the in-sample exception (no held-out split exists, so its absolute number is in-sample and fair only for shipped-vs-\texttt{v4} $\Delta$).
% [source: ledger=v5-lda-separability; eval=distill_cv_mlp; run=mri/cli.py; cache=mri/results/v5_report_regen/stats.json; config=distill_cv_5fold]
% [source: docs/analysis/2026-06-18_filter_provenance_train_heldout.md §1 (v4 train splits; bird.v1i 728 train/484 test seed 0; held-out bird 91->30, rgb_dataset_test recall 0.694->0.874); eval/build_balanced_v4_birdsplit.py]

The two layers are chosen for complementary scale coverage, not convenience. The detector's feature pyramid exposes \texttt{p3}, \texttt{p4}, and \texttt{p5} at strides $8$, $16$, and $32$: \texttt{p3} preserves the fine spatial detail that separates a small, distant drone from a bird at a few pixels, while \texttt{p5} carries the coarse semantic context (sky, clutter, gross body shape) that separates the classes at range. The intermediate \texttt{p4} is omitted because \texttt{p3} and \texttt{p5} already bracket that scale span, and the Model MRI confirms the choice rather than assuming it: its per-layer ANOVA-$F$ ranking makes \texttt{p5} the dominant discriminative block and \texttt{p3} the next, both ahead of the detection metadata (Section~\ref{sec:mri_findings}, Appendix~\ref{app:mri_report}). The pooling matches each layer's role: a $2{\times}2$ pool keeps coarse spatial structure on the high-resolution \texttt{p3}, where small targets live, while a $1{\times}1$ pool keeps only the channel response on \texttt{p5}, whose $32$-pixel cells already summarize the whole box.
% [source: classifier/mlp_verifier.py (p3 stride 8 / p4 stride 16 / p5 stride 32; P3_GRID=2x2, P5_GRID=1x1; 517-D schema); MRI per-layer ANOVA-F ranking p5 > p3 > metadata (sec:mri_findings / app:mri_report)]

Feature reuse makes the filter $46$--$72\times$ faster per detection than the patch CNN ($1.3$--$2.1$~ms vs $59$--$112$~ms; $1$--$4\%$ pipeline overhead), which is what makes \emph{per-frame} verification affordable and removes the alert-gate coupling of the predecessor design. On the shipped dataset the MLP reaches 5-fold CV $F1$ of $0.9857\pm0.0004$, rejecting $97\%$ of confuser detections while retaining $98.9\%$ of true drones. Its IR counterpart is the thermal-native \texttt{mlp\_aligned\_thermalonly} (Section~\ref{sec:ir_xmodal_verifier}).
% [source: ledger=v5-ship-per-frame,embedding-distillation-cv; run=eval/bench_speed.py]

\begin{figure}[htbp]
    \centering
    \includegraphics[width=\textwidth]{cm_filters}
    \caption{Confuser filters, detection-level held-out confusion (a detection is kept iff $P(\text{drone})\geq$ threshold). Drone side is a clean held-out test split (\texttt{rgb\_dataset}/\texttt{ir\_dset}); confuser side is out-of-distribution confuser detections. Both keep nearly all true drones while vetoing most confusers.}
    % [source: thesis_eval/eval_filter_heldout_cm.py -> thesis_eval/results/per_model_heldout/filter_heldout_cm.json (RGB mlp_v5_v4 @0.25; IR mlp_aligned_thermalonly @0.05, deploy conf 0.40)]
    \label{fig:cm_filters}
\end{figure}
